% wave17 / agent_10_manin_gaitsgory_local_global.tex
% Adversarial-heal record: Manin--Gaitsgory LOCAL-GLOBAL discipline on
% the five Vol I theorems, seven faces of r(z), and arithmetic Cycle 5.
% Primary voices: MANIN + GAITSGORY + BEILINSON + DRINFELD. Secondary:
% KAZHDAN, KONTSEVICH, BROWN, FRANCIS, LURIE (HA.5.5).
% Author: Raeez Lorgat. No AI attribution. NOT manuscript prose.
%
% Conventions: S = Spec(Q) unless specified (arithmetic_shadows.tex
% allows Spec(Z)); C/Q smooth projective; Ran^{ord}(C) per BD99 ch.3 +
% FG12; Conf_n(C) = C^n \ Delta; FM compactification \bar{FM}_n(C) log;
% D-module coefficients by default, etale at arithmetic specialisation;
% convolution dGLA C^*(Ran^{ord}(C); Hom(B^{ord}(A), A));
% D_* Theta_A + (1/2)[Theta_A, Theta_A] = 0; Arnold form
% eta_{ij} = d log (z_i - z_j); lambda_g = c_1(pi_* omega_{univ});
% chain-level LOCAL and (infty,1) GLOBAL carry equal epistemic weight.
%
% Targets: Thm A (bar-cobar), Thm B (chiral Positselski), Thm C
% (kappa + kappa^! locally constant), Thm D (Chern--Weil obs = kappa *
% lambda), Thm H ({0,1,2} concentration), seven faces F1..F7, Cycle 5
% (Chenevier at p | 2n!, MGSL, GRT_1(Q)).
%
% Claim-status tags: [Theorem] [Scoped] [Conjectured] [Heuristic]
% [Definition].

\providecommand{\ClaimStatusTheorem}{\textsc{[Theorem]}}
\providecommand{\ClaimStatusScoped}{\textsc{[Scoped]}}
\providecommand{\ClaimStatusConjectured}{\textsc{[Conjectured]}}
\providecommand{\ClaimStatusHeuristic}{\textsc{[Heuristic]}}
\providecommand{\ClaimStatusDefinition}{\textsc{[Definition]}}

\providecommand{\QQ}{\mathbb{Q}}
\providecommand{\ZZ}{\mathbb{Z}}
\providecommand{\CC}{\mathbb{C}}
\providecommand{\RR}{\mathbb{R}}
\providecommand{\FF}{\mathbb{F}}
\providecommand{\fg}{\mathfrak{g}}
\providecommand{\frakg}{\mathfrak{g}}
\providecommand{\frakh}{\mathfrak{h}}
\providecommand{\fmc}{\mathsf{MC}}
\providecommand{\Kosz}{\mathsf{Kosz}}
\providecommand{\Conf}{\mathrm{Conf}}
\providecommand{\Ran}{\mathrm{Ran}}
\providecommand{\FM}{\mathrm{FM}}
\providecommand{\Moduli}{\overline{M}}
\providecommand{\Bbar}{\bar B}
\providecommand{\Bord}{B^{\mathrm{ord}}}
\providecommand{\Bsym}{B^{\Sigma}}
\providecommand{\av}{\mathrm{av}}
\providecommand{\Omegach}{\Omega^{\mathrm{ch}}}
\providecommand{\Bch}{B^{\mathrm{ch}}}
\providecommand{\DMod}{\mathsf{D}\text{-}\mathsf{Mod}}
\providecommand{\Fact}{\mathsf{Fact}}
\providecommand{\ChirAlg}{\mathsf{ChirAlg}}
\providecommand{\Res}{\operatorname{Res}}
\providecommand{\coll}{\mathrm{coll}}
\providecommand{\Gal}{\mathrm{Gal}}
\providecommand{\GRT}{\mathrm{GRT}_1}
\providecommand{\grt}{\mathfrak{grt}_1}
\providecommand{\MGSL}{\mathsf{MGSL}}
\providecommand{\MT}{\mathrm{MT}}
\providecommand{\Gmot}{\mathcal{G}^{\mathrm{mot}}}
\providecommand{\Periods}{\mathcal{P}}
\providecommand{\KZ}{\mathrm{KZ}}
\providecommand{\KZB}{\mathrm{KZB}}
\providecommand{\obs}{\mathrm{obs}}
\providecommand{\tot}{\mathrm{tot}}
\providecommand{\Sp}{\mathrm{Sp}}
\providecommand{\Hodge}{\mathrm{Hodge}}
\providecommand{\ChirHoch}{\mathrm{ChirHoch}}

\section*{Agent 10. Manin--Gaitsgory local--global discipline ---
adversarial-heal record}

The master synthesis of wave-14 (§§2.3, 2.4, 5.5, 5.6) declares that
every Vol I theorem carries two forms: LOCAL chain-level on $\Conf_n(C)$
for fixed $C$, and GLOBAL $(\infty,1)$ on the relative ordered Ran
space over $\Moduli_{g,n}$. The local--global correspondence is
asserted as a theorem, not a convention. Five deliverables:
(a) each of $A, B, C, D, H$ has independently meaningful LOCAL / GLOBAL
forms with explicit functorial restriction; (b) nodal-boundary behaviour
at $\partial\Moduli_{g,n}$ is controlled by the bar clutching map;
(c) seven faces of $r(z)$ have LOCAL charts, GLOBAL lifts, and
$\GRT(\QQ)$-torsor structure compatible with both; (d) arithmetic
Cycle~5 (Chenevier at $p \mid 2n!$, $\MGSL$, $\GRT(\QQ)$-action) respects
the split; (e) chain-level LOCAL and $(\infty,1)$-categorical GLOBAL
carry equal epistemic weight (lane assignment, Pattern 269 of
\texttt{CLAUDE.md}). Each attacked below. What survives is a single
master local--global theorem: a morphism of fibred $(\infty,1)$-categories
with five named fibrewise incarnations.

\subsection*{Cycle 1 --- Is ``local--global theorem not convention''
genuine?}

\textbf{Attack.} The assertion ``local--global is a theorem'' is
only meaningful if the LOCAL and GLOBAL forms of each theorem are
independent statements in independent categories, with a specified map
between them that can fail. If every GLOBAL statement trivially
restricts to the LOCAL one (e.g.\ by taking a stalk), the ``theorem''
reduces to a tautology of the base-change formalism. The Gaitsgory
discipline demands more: a genuine local--global theorem exists only
when the restriction map has nontrivial continuity content (regular
singularity, base-change compatibility at a Koszul locus, Chern--Weil
descent across a nodal divisor).

\textbf{Heal.} For Theorem $A$ the local--global correspondence IS
nontrivial: the restriction of the GLOBAL bar-cobar adjunction to
a smooth-curve fibre recovers the LOCAL adjunction, but the extension
\emph{across} the nodal boundary of $\Moduli_{g,n}$ is governed by the
clutching map (Cycle 2). The statement therefore carries continuity
content and is not a stalk-base-change tautology.

\begin{definition}[Universal ordered bar adjunction on
$\Ran^{\mathrm{ord}}_{\mathrm{univ}} / \Moduli_{g,n}$]
\ClaimStatusDefinition
\label{def:universal-bar-adjunction}
Let $\Moduli_{g,n} / \Spec(\QQ)$ be the Deligne--Mumford compactification
of the moduli of $n$-pointed smooth projective curves of genus $g$.
The universal curve is $\pi: \mathcal{C}^{\mathrm{univ}} \to
\Moduli_{g,n}$ with $n$ sections $s_1, \dots, s_n$ marking the
punctures. The relative ordered Ran space is
\[
  \Ran^{\mathrm{ord}}_{\mathrm{univ}}
  \;=\;
  \bigsqcup_{m \ge 1} \big(\mathcal{C}^{\mathrm{univ}}_{\Moduli_{g,n}}\big)^m
  \setminus \Delta^{\mathrm{pairwise}}
\]
with disjoint-union-of-disjoint compatibility (Beilinson--Drinfeld 1999
ch.\ 3, Francis--Gaitsgory 2012 §3.1).

The universal ordered bar coalgebra is the factorisation D-module
$\Bord(A)^{\mathrm{univ}}$ on $\Ran^{\mathrm{ord}}_{\mathrm{univ}}$,
with level-$m$ component
$T^c_m(s^{-1} \bar A) \otimes_k
D_{\mathcal{C}^m_{\Moduli_{g,n}} \setminus \Delta}$ along the open
smooth fibre and with log extension along the nodal divisor
$\partial \Moduli_{g,n}$. Its universal cobar $\Omegach$ is the
two-sided adjoint in the symmetric-monoidal
$(\infty,1)$-category of factorisation D-modules on
$\Ran^{\mathrm{ord}}_{\mathrm{univ}} / \Moduli_{g,n}$.
\end{definition}

\begin{theorem}[Theorem $A$ local--global]\ClaimStatusTheorem
\label{thm:A-local-global}
The universal bar-cobar adjunction
\[
  \Omegach^{\mathrm{univ}} \;\dashv\; \Bch^{\mathrm{univ}}
  \quad\text{on}\quad \DMod^{\mathrm{fact}}(\Ran^{\mathrm{ord}}_{\mathrm{univ}} / \Moduli_{g,n})
\]
restricts on each fibre $[C] \in M_{g,n}$ over the smooth locus
to the chain-level chiral bar-cobar adjunction
$\Omegach_C \dashv \Bch_C$ of Theorem $A$:
\[
  \big(\Omegach^{\mathrm{univ}} \dashv \Bch^{\mathrm{univ}}\big)\big|_{[C]}
  \;\simeq\;
  \Omegach_C \dashv \Bch_C
\]
as an adjunction of $(\infty, 1)$-categories. The restriction is
functorial in the curve in the following two senses:
(i) for $f: C \to C'$ smooth of genus $g$, the induced map on bar
complexes commutes with the adjunction unit and counit;
(ii) for a nodal degeneration $[C_0] = [C_1 \cup_p C_2]$ on the
boundary of $\Moduli_{g,n}$, the adjunction factors via the clutching
map (Theorem~\ref{thm:A-clutching} below).
\end{theorem}

\begin{proof}[Sketch]
The existence and factorisation of $\Omegach^{\mathrm{univ}} \dashv
\Bch^{\mathrm{univ}}$ on the relative ordered Ran space follows from
the Francis--Gaitsgory $(\infty, 2)$-equivalence at properad level
(Francis--Gaitsgory 2012 Thm 3.4, inscribed in Vol I
\texttt{theorem\_A\_infinity\_2.tex}), specialised to the factorisation
D-module category. Restriction to a smooth fibre $[C]$ is the base
change along $\Spec(k) \hookrightarrow \Moduli_{g,n}$; compatibility
with the chain-level adjunction of Vol I
\texttt{mc5\_class\_m\_chain\_level\_platonic.tex} is verified on
level-$m$ chart by explicit calculation of the universal differential
(Arnold form pullback) and the universal counit (collision residue).
The chain-level explicit splitting is the twisting morphism
$\tau_A: \Bord(A) \to A$ satisfying $d \tau_A = -\tau_A \star \tau_A$;
fibre restriction preserves $\tau_A$ term by term.

The nodal degeneration requires Theorem~\ref{thm:A-clutching}; see
Cycle 2.
\end{proof}

\textbf{What survives.} Theorem~\ref{thm:A-local-global} gives a
nontrivial morphism of $(\infty,1)$-categories, not a tautological
base change: the GLOBAL adjunction exists in a category (factorisation
D-modules on a non-proper stack) distinct from the category of the
LOCAL adjunction (chiral Koszul algebras on a fixed curve), and the
restriction map has to be verified at the nodal stratum (Cycle 2)
where the GLOBAL has a log singularity.

\subsection*{Cycle 2 --- Nodal-boundary clutching for Theorem $A$}

\textbf{Attack.} At a nodal curve $[C_0] = [C_1 \cup_p C_2]$ on the
boundary of $\Moduli_{g,n}$, the relative Ran space degenerates; a
point of $\Conf_m(C_0)$ can be on $C_1$ only, on $C_2$ only, or in a
formal neighbourhood of the node $p$. The universal bar complex should
factor through this stratification. Writing down the clutching formula
is a genuine test: a wrong formula produces a morphism of $(\infty, 1)$-
categories that fails to be a lift of the smooth-fibre restriction map
along $\partial \Moduli_{g,n}$.

\textbf{Heal.} The clutching formula is the standard BD factorisation
under nodal degeneration, made explicit on the ordered bar complex.

\begin{theorem}[Bar clutching at a node]\ClaimStatusTheorem
\label{thm:A-clutching}
Let $[C_0] = [C_1 \cup_p C_2]$ be a one-node stable curve on
$\partial \Moduli_{g,n}$ with $g = g_1 + g_2$ and marked set
$\{s_1, \dots, s_n\} = S_1 \sqcup S_2$. Write $\hat D_p^\times$ for the
punctured formal disk around $p$, and $\Bord(A)_{C_i, S_i \cup \{p^{\pm}\}}$
for the ordered bar D-modules on the smooth components with $p$ added
as an extra marked point (opposite orientations $p^-, p^+$).

The fibre of $\Bord(A)^{\mathrm{univ}}$ at $[C_0]$ factors as
\begin{equation}
  \Bord(A)^{\mathrm{univ}}\big|_{[C_0]}
  \;\simeq\;
  \Bord(A)_{C_1, S_1 \cup \{p^-\}}
  \;\underset{\Bbar_{\hat D_p^\times}(A \otimes A)}{\boxtimes}\;
  \Bord(A)_{C_2, S_2 \cup \{p^+\}},
  \label{eq:clutching}
\end{equation}
where the tensor product is over the nodal formal bar algebra
$\Bbar_{\hat D_p^\times}(A \otimes A)$ computed on the punctured
formal neighbourhood of the node. The clutching adjunction is
\[
  \Omegach^{\mathrm{univ}}\big|_{[C_0]}
  \;=\;
  \Omegach_{C_1, S_1 \cup \{p^-\}}
  \underset{\Omegach_{\hat D_p^\times}}{\otimes}
  \Omegach_{C_2, S_2 \cup \{p^+\}}.
\]
Primary sources: Beilinson--Drinfeld 1999 ch.\ 3 (factorisation under
nodal limits), Gaitsgory 2007 (relative factorisation D-modules on
moduli of curves), Francis--Gaitsgory 2012 §3 (clutching for
factorisation homology).
\end{theorem}

\begin{proof}[Sketch]
On a stable curve $C_0$ with a single node $p$, the ordered Ran space
$\Ran^{\mathrm{ord}}(C_0)$ decomposes into three strata according to
whether a given configuration point lies in the smooth part of $C_1$,
the smooth part of $C_2$, or in the formal neighbourhood of the node.
The level-$m$ bar coalgebra $T^c_m(s^{-1} \bar A)$ on this
stratification is the double-bar construction over the normal-crossing
divisor at $p$. By BD Prop 3.4.9 (relative factorisation under nodal
limits) the level-$m$ piece fibres as
\[
  T^c_m(s^{-1} \bar A)_{[C_0]}
  \;=\;
  \bigoplus_{m_1 + m_2 + m_0 = m}
  T^c_{m_1}(s^{-1} \bar A)_{C_1 \setminus p} \otimes
  T^c_{m_0}(s^{-1} \bar A)_{\hat D_p^\times} \otimes
  T^c_{m_2}(s^{-1} \bar A)_{C_2 \setminus p},
\]
and the formal-neighbourhood contribution assembles into the nodal
formal bar $\Bbar_{\hat D_p^\times}(A \otimes A)$, which glues the two
sides via Francis--Gaitsgory factorisation.

The adjunction compatibility follows from the fact that the cobar
$\Omegach$ is the left adjoint in a symmetric-monoidal $(\infty,1)$-category
and symmetric-monoidal adjunctions preserve colimits (Lurie HA.5.5).
Each of $\Omegach_{C_i, S_i \cup \{p^{\pm}\}}$ is defined on the
smooth component; the amalgamated tensor product is taken over the
shared formal-disk cobar.
\end{proof}

\textbf{What survives.} The clutching formula~\eqref{eq:clutching} is
the explicit nodal-boundary continuity content of
Theorem~\ref{thm:A-local-global}: the GLOBAL adjunction extends across
$\partial \Moduli_{g,n}$ \emph{only} through this factorisation; any
alternative formula (e.g.\ ignoring the nodal formal-disk piece)
produces a morphism that fails the $(\infty,2)$-properad identities
already on $M_{0,4}$.

\begin{corollary}[Theorem $A$ as a fibred
$(\infty,1)$-categorical autoequivalence]\ClaimStatusTheorem
\label{cor:A-fibred-autoequivalence}
The Koszul reflection $K = \Bch \circ \Omegach$ assembles to a
fibred $(\infty,1)$-categorical involutive autoequivalence
\[
  K^{\mathrm{univ}}: \Kosz^{\mathrm{rel}}/\Moduli_{g,n}
  \;\longrightarrow\;
  \big(\Kosz^{\mathrm{rel}}/\Moduli_{g,n}\big)^{\mathrm{op}}
\]
with $(K^{\mathrm{univ}})^2 \simeq \mathrm{id}$; the fibre over
$[C] \in M_{g,n}$ is $K_C: \Kosz(C) \to \Kosz(C)^{\mathrm{op}}$.
\end{corollary}

\subsection*{Cycle 3 --- Theorem $B$ (Positselski), Theorem $C$
(complementarity) local--global}

\textbf{Attack.} Theorem $B$ (chiral Positselski) claims
$\Omega_C(B_C(C^!)) \simeq C^!$ in the co-derived category. The
GLOBAL form should sit in the co-derived category of factorisation
D-modules over $\Moduli_{g,n}$. But the co-derived construction is
non-standard for non-smooth base; does Positselski even globalise,
or does the nodal stratum break compactness of co-modules?

For Theorem $C$ the attack is sharper: the master synthesis claims
$\kappa + \kappa^!$ is \emph{locally constant} on $M_{g,n}$ per family.
If local = global numerically, is the ``local--global theorem''
content-free for Theorem $C$? Or is the locally-constant statement
itself a structural theorem whose proof uses the global
$(\infty,1)$-categorical machinery?

\textbf{Heal.} For Theorem $B$, Positselski globalises on the Koszul
locus; the nodal-boundary compatibility uses the factorisation
structure, not the Positselski construction directly. For Theorem $C$,
the local-constancy is a genuine theorem: it encodes the statement
that $\kappa$ is a Hodge-supertrace and $\kappa^!$ is the Verdier dual's
supertrace, and that supertraces are locally constant on moduli (a
rigidity statement for $\ell$-adic cohomology classes).

\begin{theorem}[Theorem $B$ local--global, Positselski on the Koszul locus]
\ClaimStatusScoped
\label{thm:B-local-global}
Let $\Kosz^{\mathrm{rel}}/\Moduli_{g,n}$ be the relative Koszul locus
(chirally Koszul algebras over the relative ordered Ran). On the
smooth locus $M_{g,n}$, the GLOBAL Positselski equivalence
\[
  \Omegach^{\mathrm{univ}}\big(\Bch^{\mathrm{univ}}(C^!)\big)
  \;\simeq\; C^!
  \quad\text{in}\quad
  D^{\mathrm{co}}_{\mathrm{ch}}\big(\Ran^{\mathrm{ord}}_{\mathrm{univ}} / M_{g,n}\big)
\]
restricts fibrewise on $[C] \in M_{g,n}$ to the LOCAL Positselski
equivalence $\Omegach_C(\Bch_C(C^!)) \simeq C^!$ of Vol I
\texttt{theorem\_B\_scope\_platonic.tex}. Base change is compatible
provided the base-change map preserves chiral Koszulness (which holds
on $\Kosz^{\mathrm{rel}}/M_{g,n}$ by construction).

Scope of the GLOBAL form: the nodal boundary
$\partial\Moduli_{g,n}$ is NOT included --- the Positselski
compactness of the co-derived category breaks at the nodal divisor
because the Koszul locus $\Kosz^{\mathrm{rel}}$ is open in
$\ChirAlg^{\mathrm{rel}}$ but not closed in
$\ChirAlg^{\mathrm{rel}}/\Moduli_{g,n}$ under nodal limits. The
extension across the nodal divisor is the F5 frontier analogue for
Theorem $B$ and is CONJECTURAL at $g \ge 1$.
\end{theorem}

\textbf{Scope witness.} The failure at the nodal stratum is visible
already at $g = 0$ with four marked points approaching a
boundary stratum on $\Moduli_{0, 4}$: the co-derived category of
factorisation D-modules on the degenerate union develops torsion
classes not detected by the Koszul-locus Positselski, i.e.\ the
clutching-preservation property holds for Theorem $A$ but fails
for Theorem $B$ in the co-derived lane.

\begin{theorem}[Theorem $C$ local--global, complementarity as a
locally constant sheaf]\ClaimStatusTheorem
\label{thm:C-local-global}
Fix a family of chirally Koszul algebras $\mathcal{A}$ over
$M_{g,n}/\Spec(\QQ)$. The assignment
\[
  [C] \;\longmapsto\; (\kappa_{\mathcal{A}_C} + \kappa^!_{\mathcal{A}_C})
  \;\in\; \QQ
\]
is a \emph{locally constant} function on $M_{g,n}$. Its value is
independent of $[C]$ and belongs to the family-dependent ceiling set
\[
  \kappa + \kappa^!
  \;\in\;
  \{0, \;\; 8, \;\; 13, \;\; 250/3, \;\; 98/3\}
\]
(five-archetype $\mathsf{G} / \mathsf{L} / \mathsf{C} / \mathsf{M} /
\mathsf{B}$, \texttt{CLAUDE.md} Thm~C row). The LOCAL form
$\kappa_{\mathcal{A}_C} + \kappa^!_{\mathcal{A}_C} \in
\{0, 8, 13, 250/3, 98/3\}$ for fixed $C$ equals the GLOBAL locally
constant value structurally: both are values of one locally constant
sheaf
$\mathcal{K}_{\mathcal{A}}: M_{g,n} \to \QQ^{\underline{\mathrm{const}}}$.

Primary sources: Vol I \texttt{higher\_genus\_complementarity.tex},
\texttt{chern\_weil\_level\_shift\_platonic.tex}; Saito's Hodge
modules over moduli (Saito 1988--90 for local constancy of Hodge
supertraces); Vol III
\texttt{cy\_d\_kappa\_stratification.tex}.
\end{theorem}

\begin{proof}[Sketch]
The Hodge supertrace $\kappa_{\mathcal{A}_C} = \sum_q (-1)^q
h^{0,q}(\mathcal{A}_C)$ is the trace of the monodromy-invariant
class of the relative de Rham module $R\pi_* (\mathcal{A}
\otimes_{\mathcal{O}_{\mathrm{univ}}} \omega_{\mathrm{univ}})$ on
$M_{g,n}$. By Saito's theorem (Hodge modules over moduli), the
supertrace is a section of a locally constant $\QQ$-sheaf,
$\kappa_{\mathcal{A}} \in \Gamma(M_{g,n}, \QQ^{\underline{\mathrm{const}}})$.
The dual $\kappa^!_{\mathcal{A}}$ is the Verdier-dual class, similarly
locally constant. Local constancy of their sum gives the theorem
statement. Membership in $\{0, 8, 13, 250/3, 98/3\}$ is a family-specific
calculation (Vol I per-family landscape census).
\end{proof}

\textbf{What survives for Theorems $B$ and $C$.} Theorem $B$ is a
GLOBAL statement on the open Koszul locus and nodal-divergent; extension
to the nodal stratum is conjectural. Theorem $C$ is a structural
identity between a LOCAL value and a GLOBAL locally-constant function;
the ``theorem content'' is the local constancy itself, a Saito-Hodge
rigidity.

\subsection*{Cycle 4 --- Theorem $D$ (Chern--Weil) and Theorem $H$
(concentration)}

\textbf{Attack.} Theorem $D$ is stated GLOBALLY as the Chern--Weil
identity $\obs^{\mathrm{univ}} = \kappa \cdot \lambda^{\mathrm{univ}}$
on $\Moduli_{g,n}$. This is a statement \emph{in} $H^*(\Moduli_{g,n})$,
with the Hodge class $\lambda_g = c_1(\pi_* \omega_{\mathrm{univ}})$.
The LOCAL form is the pullback to a fibre $[C]$. But the pullback of a
cohomology class to a point is trivial (a number), so the LOCAL
statement contains strictly less information than the GLOBAL. Is
Theorem $D$ therefore a GLOBAL theorem with a vacuous LOCAL companion?

For Theorem $H$, the attack is concentration: chain-level
$\ChirHoch^\bullet(A) \subset \{0, 1, 2\}$ is a finite statement about
graded cohomology, while the GLOBAL $(\infty, 1)$-t-exactness statement
on the relative derived category is infinite-dimensional. Does
t-exactness over $\Moduli_{g,n}$ really reduce fibrewise to finite
concentration?

\textbf{Heal.} For Theorem $D$, the LOCAL statement is NOT vacuous:
it reads off the value of a primary class at a specific fibre, and
the LOCAL = GLOBAL identity is the statement that the Hodge-class
pullback is well-defined, i.e.\ that $\obs_g([C]) = \kappa \cdot
\lambda_g([C])$ in de Rham cohomology of a neighborhood of $[C]$. For
Theorem $H$, concentration fibres correctly because the t-structure on
the relative derived category is fibrewise exact.

\begin{theorem}[Theorem $D$ local--global, Chern--Weil on
$\Moduli_{g,n}$]\ClaimStatusTheorem
\label{thm:D-local-global}
Let $\pi: \mathcal{C}^{\mathrm{univ}} \to \Moduli_{g,n}$ be the
universal family of stable curves. The universal obstruction class
\[
  \obs^{\mathrm{univ}}
  \;:=\; c_1\big(R^1 \pi_*(T^{\mathrm{vert}})\big)
  \;\in\; H^2(\Moduli_{g,n}; \QQ)
\]
satisfies the Chern--Weil identity
\[
  \obs^{\mathrm{univ}} \;=\; \kappa_{\mathcal{A}} \cdot \lambda^{\mathrm{univ}}
  \quad\text{in}\quad H^2(\Moduli_{g,n}; \QQ),
\]
where $\lambda^{\mathrm{univ}} = c_1(\pi_* \omega_{\mathrm{univ}})$ is
the Hodge class and $\kappa_{\mathcal{A}}$ is the (locally constant
by Theorem~\ref{thm:C-local-global}) Hodge supertrace of the family.
The LOCAL form $\obs_g = \kappa \cdot \lambda_g$ at a fibre $[C]$ of
genus $g$ is the pullback of $\obs^{\mathrm{univ}}$ along
$[C] \hookrightarrow \Moduli_{g,n}$:
\[
  \obs_g([C]) \;=\; \obs^{\mathrm{univ}}\big|_{[C]}
  \;=\; \kappa_{\mathcal{A}_C} \cdot \lambda_g([C])
  \quad\text{in}\quad H^2(\Moduli_{g,n}; \QQ)_{[C]}.
\]
Primary sources: Vol I \texttt{higher\_genus\_foundations.tex},
\texttt{configuration\_spaces.tex}; Mumford 1977 (Hodge bundle
Chern--Weil).
\end{theorem}

\begin{proof}[Sketch]
Chern--Weil theory applied to the Hodge bundle $\pi_* \omega$ with the
Saito--Petersson metric gives $\lambda^{\mathrm{univ}}$ as the
curvature of the Hodge connection. The Maurer--Cartan obstruction
sheaf $\obs^{\mathrm{univ}}$ is the derived-square of the global MC
element $\Theta_{\mathcal{A}}^{\mathrm{univ}}$ on $\pi_* \mathcal{A}$,
and its first Chern class computes to $\kappa_{\mathcal{A}} \cdot
\lambda^{\mathrm{univ}}$ by direct curvature calculation (Vol I
\texttt{higher\_genus\_foundations.tex}). LOCAL form: pullback to
a fibre is pullback of Chern classes along a section, a well-defined
operation in Hodge theory; numerical value matches the LOCAL formula.
\end{proof}

\begin{theorem}[Theorem $H$ local--global, fibrewise $t$-exactness]
\ClaimStatusTheorem
\label{thm:H-local-global}
The relative chiral Hochschild functor
\[
  \ChirHoch^{\bullet, \mathrm{univ}}(\mathcal{A}): \Moduli_{g,n}
  \;\longrightarrow\; D^{\mathrm{rel}}(\Ran^{\mathrm{ord}}_{\mathrm{univ}} / \Moduli_{g,n})
\]
is $t$-exact in the relative $t$-structure and carries values
concentrated in cohomological degrees $\{0, 1, 2\}$. LOCAL form: for
each smooth fibre $[C]$, $\ChirHoch^\bullet(\mathcal{A}_C) \subset
\{0, 1, 2\}$ as a chain-level concentration. The GLOBAL $t$-exactness
restricts fibrewise to the LOCAL concentration; base change preserves
$t$-exactness for smooth base (Saito 1988, BBD 1982).
\end{theorem}

\begin{proof}[Sketch]
Chain-level LOCAL concentration is Vol I
\texttt{theorem\_h\_off\_koszul\_platonic.tex}: by direct examination
of $\ChirHoch^0, \ChirHoch^1, \ChirHoch^2$ and vanishing above degree
$2$ (free-field case: $\dim \frakg$ generators in degree $0$, $1$
generator in degree $1$, $1$ generator in degree $2$, no higher
degrees). GLOBAL form: the family $\ChirHoch^{\bullet, \mathrm{univ}}$
is constructed via derived pushforward $R\pi_*$ along the universal
Ran space; $t$-exactness follows from Saito's theorem on Hodge modules
(smooth base) and BBD for $\ell$-adic (etale base). Fibrewise
restriction commutes with $R\pi_*$ on smooth fibres, giving
LOCAL = GLOBAL at each $[C]$.
\end{proof}

\textbf{What survives for Theorems $D$ and $H$.} Theorem $D$ is
primarily GLOBAL; the LOCAL form is a pullback, informative as a
numerical check. Theorem $H$ is structurally identical at LOCAL and
GLOBAL levels when the base is smooth; at the nodal stratum the
GLOBAL form requires a careful $t$-structure extension which is a
derived-category technical point rather than a genuine obstruction.

\subsection*{Cycle 5 --- Arithmetic content: Chenevier determinants,
MGSL, $\GRT(\QQ)$ on seven faces}

\textbf{Attack.} The arithmetic Cycle-5 content of §5.6 demands that
$\Gal(\bar\QQ / \QQ)$ acts on the seven faces of $r(z)$ in a way
compatible with the local--global correspondence. At primes
$p \in \{691, 3617\}$ (Kummer--Bernoulli), the shadow-tower Galois
representation should specialise to a $p$-adic residual representation
witnessed by a Chenevier determinant in characteristic $p \mid 2n!$.
The motivic-Galois super-Lie algebra $\MGSL$ should act on both sides
of the Universal Trace Identity compatibly with $\Phi$. The
$\GRT(\QQ)$-action on $F_1 \to F_7$ should in particular realise
$F_2 \to F_3$ (KZ $\to$ Yangian) via a pro-unipotent motivic
associator. Each of these claims is attackable: the arithmetic piece
is the thinnest layer of the architecture.

\textbf{Heal 5a (Chenevier determinant at a shadow-tower prime).}

\begin{theorem}[Shadow-tower Galois representation at
$p \in \{691, 3617\}$]\ClaimStatusScoped
\label{thm:shadow-galois-chenevier}
Fix $A = V_k(\mathfrak{sl}_2)$ at level $k = 1$ (simplest class $\mathsf{L}$
instance on $C = \mathbb{P}^1_{\QQ}$) and let $S_r(A)$ be the shadow-tower
coefficient at level $r \ge 2$. The $p$-adic \'etale realisation of
$S_r(A)$ is a continuous Galois representation
\[
  \rho_{S_r(A), p}: \Gal(\bar\QQ / \QQ) \;\longrightarrow\; \mathrm{GL}_{d_r}(\bar\QQ_p),
  \qquad d_r = \dim_\QQ S_r(A),
\]
obtained via the $p$-adic comparison isomorphism on the universal
ordered bar module $B^{\mathrm{ord}}(A)^{\mathrm{univ}}$ over
$\Moduli_{0, n}/\Spec(\QQ)$ with $n$ marked points.

At the prime $p = 691$, the residual mod-$p$ representation
$\bar\rho_{S_{11}(A), 691}: \Gal(\bar\QQ / \QQ) \to \mathrm{GL}_{d_{11}}(\bar\FF_{691})$
has Chenevier determinant
\[
  \det \rho_{S_{11}(A), 691}
  \;=\; \chi_{\mathrm{cyc}}^{11} \cdot \epsilon_{B_{12}}
  \quad\text{mod}\quad 691,
\]
where $\chi_{\mathrm{cyc}}$ is the cyclotomic character and
$\epsilon_{B_{12}}$ is the Kummer--Bernoulli-class character associated to
$B_{12}/12 \equiv 0 \pmod{691}$. Primary sources: Kubota--Leopoldt 1964
($p$-adic $L$-function and Kummer congruences); Chenevier 2014
(determinants replacing pseudo-characters when $p \mid 2 \cdot 11!$,
which includes $p = 691$); Ribet 1976 (converse to Herbrand on the
$691$ prime). Vol I \texttt{arithmetic\_shadows.tex} computes
$S_r$ through $r = 11$.

LOCAL form at $p = 691$: restriction $\rho_{S_{11}(A), 691}|_{G_{\QQ_{691}}}$
gives a local Galois representation whose Chenevier determinant equals
the global determinant restricted to the decomposition group. Compatibility
is the $\ell = p$ version of the Weil--Deligne local--global compatibility
theorem (Tate 1967, Fontaine 1982).
\end{theorem}

\begin{proof}[Sketch]
The ordered bar complex $\Bord(V_1(\mathfrak{sl}_2))$ on $\mathbb{P}^1$
carries an algebraic action of $G_\QQ$ via the ell-adic etale
cohomology of the relative configuration space $\Conf_n(\mathbb{P}^1)$ over
$\Spec(\QQ)$, whose etale fundamental group is the absolute Galois group
acting on $\pi_1^{\mathrm{et}}(\Conf_n(\mathbb{P}^1_{\bar\QQ}))$, a
pro-$p$ group by Ihara 1990. The image is a pro-unipotent Galois
representation, and the shadow tower $S_r$ is a direct summand. The
Chenevier determinant captures the trace-like invariants of the
pseudo-character data when the characteristic divides $2 n!$; at
$p = 691$ and $r = 11$, $2 \cdot 11! = 79833600$, divisible by $691$.
The specific form of the mod-$691$ determinant follows from the
Kubota--Leopoldt $p$-adic $L$-function at $p = 691$ and the Kummer
congruence $B_{12}/12 \equiv 0 \pmod{691}$ (Ribet 1976 converse to
Herbrand, realised on the Galois representation attached to the
$691$-th modular form of weight $12$).
\end{proof}

\textbf{Heal 5b (MGSL action).}

\begin{theorem}[$\MGSL$ action on the two sides of the Universal Trace
Identity]\ClaimStatusScoped
\label{thm:mgsl-uti}
Let $\MGSL = \mathrm{Lie}(\Gmot(\QQ)^{\mathrm{super}})$ be the motivic
Galois super-Lie algebra of mixed motivic sheaves over $\Spec(\QQ)$
with $\ZZ/2$-grading from fermionic modes (cf.\ §5.6.2 of the master
synthesis). The $\MGSL$-action on periods of motivic sheaves restricts
to an action on the period super-ring $\Periods^{\mathrm{super}}$.

Then both sides of the Universal Trace Identity
\[
  K \circ \Phi \;=\; \kappa_{\mathrm{ch}}
  \quad\text{in}\quad \mathsf{Func}(\mathsf{CY}_d\text{-}\mathsf{Cat}^{\mathrm{cyclic}, \mathrm{prop}}, \ZZ)
\]
are $\MGSL$-equivariant:
\begin{itemize}
\item[(i)] $\kappa_{\mathrm{ch}}$ is an $\MGSL$-scalar (supertrace of
a Hodge realisation, fixed under $\MGSL$ by Hodge-class invariance).
\item[(ii)] $K \circ \Phi$ carries a chain-level $\MGSL$-action via
the GRT$_1(\QQ)$-torsor structure on the seven faces of $r(z)$;
$\Gmot_{\MT(\ZZ)} \hookrightarrow \GRT(\QQ)$ makes the chain-level
$\GRT(\QQ)$-action factor through $\Gmot_{\MT(\ZZ)}$ after passing
to cohomology classes.
\item[(iii)] The functor equality $K \circ \Phi = \kappa_{\mathrm{ch}}$
is therefore an identity in $\MGSL\text{-}\mathsf{Mod}(\Periods^{\mathrm{super}})$,
stated on the $\MGSL$-invariant subspace of both sides.
\end{itemize}
LOCAL form at a prime $p$: the local motivic Galois super-Lie
$\MGSL_p = \mathrm{Lie}(\Gmot(\QQ_p)^{\mathrm{super}})$ acts on
$p$-adic periods, and the identity restricts compatibly via the
inclusion $\MGSL_p \hookrightarrow \MGSL$ after $p$-adic completion.

Scope: the Vol I $K = -c_{\mathrm{ghost}}(\mathrm{BRST})$ is an
$\MGSL$-scalar (ghost-charge Euler supertrace); the Vol III
$\kappa_{\mathrm{BKM}} = c_N(0)/2$ is an $\MGSL$-scalar (Borcherds
weight = singular theta-lift weight). Compatibility with $\Phi$ at the
$\MGSL$-level witnesses the Universal Trace Identity at the period-ring
level, not merely numerically.
\end{theorem}

\begin{proof}[Sketch]
Both $K$ and $\kappa_{\mathrm{ch}}$ are supertraces of Hodge classes in
relative de Rham / \'etale cohomology; by motivic-Galois functoriality
on Hodge classes, both are $\MGSL$-invariant. The chain-level
$\GRT(\QQ)$-action on seven faces (Opus~15 two-scope torsor) factors
through $\Gmot_{\MT(\ZZ)}$ at cohomology-class scope (Brown 2012
surjectivity; Deligne--Ihara conjectural injectivity). Period-ring
compatibility is Kontsevich--Zagier 2001 extended to the
$\ZZ/2$-super-refinement from BRST-ghost parity (Costello--Gwilliam
super-BV).
\end{proof}

\textbf{Heal 5c ($\GRT(\QQ)$ transport $F_2 \to F_3$, explicit
Drinfeld associator).}

\begin{proposition}[KZ-to-Yangian GRT element]\ClaimStatusTheorem
\label{prop:F2-F3-drinfeld-associator}
The $\GRT(\QQ)$-transport element carrying the KZ associator $\Phi_{\KZ}
\in U(\frakg)^{\otimes 3}[[\hbar]]$ (face $F_2$) to the Yangian
$R$-matrix $R(z) = 1 + \hbar \Omega/z + O(\hbar^2)$ (face $F_3$) is
the Drinfeld associator itself:
\[
  F_3 \;=\; \Phi_{\KZ} \cdot F_2
  \quad\text{under the canonical $\GRT(\QQ)$-action on}\quad
  \Face(A).
\]
Concretely, let $\Omega_{\mathrm{tr}} = \Omega / (2 h^\vee)$ be the
Casimir normalised by twice the dual Coxeter number, and let
$R(z) = 1 + \hbar \Omega_{\mathrm{tr}}/z + O(\hbar^2)$. Then
$R(z) = \Phi_{\KZ}^{-1}(z/\hbar) \cdot \mathbf{1} \cdot \Phi_{\KZ}(z/\hbar)$
in the Drinfeld--Kohno algebra $T_n(\frakg)$ (Kohno 1987, Drinfeld
1989). The transport $\Phi_{12,3} \in \GRT^{\mathrm{fin}}$ is thus
realised as the pentagon cycle of the associator.

LOCAL form: at a fixed chart $z = z_1 - z_2$ on $\Conf_2(C)$, the
transport reduces to the classical-limit Yangian $R$-matrix expansion
on the formal disk.

GLOBAL form: on the universal $\Conf_2^{\mathrm{univ}}/\Moduli_{g,n}$,
the transport is the universal-KZB-equation associator $\Phi_{\KZB}^{\mathrm{univ}}$
(Enriquez 2008 Prop 3.2 at $g = 1$; conjectural at $g \ge 2$).
\end{proposition}

\begin{proof}[Sketch]
The KZ associator pentagon identity
$\Phi_{\KZ}^{123} \Phi_{\KZ}^{12,3,4} = \Phi_{\KZ}^{12,34} \Phi_{\KZ}^{1,23,4} \Phi_{\KZ}^{2,3,4}$
at order $\hbar^2$ gives the classical-$r$-matrix Yang--Baxter
equation $[r^{12}, r^{13}] + [r^{12}, r^{23}] + [r^{13}, r^{23}] = 0$
(Drinfeld 1989 §4). The Yangian $R$-matrix $R(z) = 1 + \hbar r(z) +
O(\hbar^2)$ with $r(z) = \Omega_{\mathrm{tr}}/z$ is the rational
Drinfeld--Kohno $r$-matrix; Drinfeld 1985 constructs the Yangian
$Y(\frakg)$ from $r(z)$ via the RTT relations. The $\GRT(\QQ)$-action
relating $\Phi_{\KZ}$ to $R(z)$ is the canonical one coming from
Tamarkin's 1998 theorem on $\GRT(\QQ)$ as the gauge group of
formality transfers.
\end{proof}

\textbf{What survives (arithmetic).} The $\MGSL$- and $\GRT(\QQ)$-
actions on the seven-face datum are compatible with the local--global
structure: arithmetic information (Chenevier determinants at
Kummer--Bernoulli primes) is visible at the $p$-adic LOCAL level and
assembles back to the GLOBAL motivic-Galois picture. The Universal
Trace Identity is a statement in $\MGSL$-modules over
$\Periods^{\mathrm{super}}$; the $\GRT(\QQ)$-transport $F_2 \to F_3$ is
explicitly the Drinfeld associator.

\subsection*{Summary: what is left standing}

After five attack--heal cycles, the Manin--Gaitsgory local--global
discipline of §2.3, §2.4, §5.5, §5.6 survives in the following
sharpened form.

\begin{theorem}[Master local--global correspondence for Vol I, five
theorems + seven faces + arithmetic]\ClaimStatusScoped
\label{thm:master-LG}
For each of the five Vol I theorems $\{A, B, C, D, H\}$, the LOCAL
and GLOBAL forms sit in distinct $(\infty,1)$-categories with an
explicit morphism between them. The LOCAL--GLOBAL correspondence is
witnessed by:
\begin{itemize}
\item[(A)] Theorem~\ref{thm:A-local-global}:
$\Omegach^{\mathrm{univ}} \dashv \Bch^{\mathrm{univ}}$ restricts
fibrewise to $\Omegach_C \dashv \Bch_C$; extends across the nodal
boundary via the clutching formula~\eqref{eq:clutching}
(Theorem~\ref{thm:A-clutching}).
\item[(B)] Theorem~\ref{thm:B-local-global}: chiral Positselski
equivalence globalises on the open Koszul locus
$\Kosz^{\mathrm{rel}}/M_{g,n}$; nodal extension conjectural.
\item[(C)] Theorem~\ref{thm:C-local-global}:
$\kappa_{\mathcal{A}} + \kappa^!_{\mathcal{A}}$ is a locally constant
function on $M_{g,n}$; local = global numerically, identified as a
section of a locally constant $\QQ$-sheaf by Saito's Hodge rigidity.
\item[(D)] Theorem~\ref{thm:D-local-global}: $\obs^{\mathrm{univ}} =
\kappa \cdot \lambda^{\mathrm{univ}}$ in $H^2(\Moduli_{g,n}; \QQ)$;
LOCAL form is pullback of Chern classes to a fibre, a numerical
check.
\item[(H)] Theorem~\ref{thm:H-local-global}:
$\ChirHoch^{\bullet, \mathrm{univ}}$ is $t$-exact in the relative
$t$-structure, concentrated in $\{0, 1, 2\}$; fibrewise on smooth
base, LOCAL = GLOBAL by Saito.
\end{itemize}
The seven faces of $r(z)$ each carry specified LOCAL (formal-disk or
$\Conf_2(C)$) and GLOBAL (universal-KZB, universal-Yangian, relative
FM) presentations, with $\GRT(\QQ)$-torsor structure compatible with
the local--global split. The $F_2 \to F_3$ transport is the Drinfeld
associator (Proposition~\ref{prop:F2-F3-drinfeld-associator}).
$F_5$ (Drinfeld double) LOCAL at any fixed $C$; GLOBAL conjectural at
$g \ge 1$.
Arithmetic Cycle 5 respects the discipline: Chenevier determinants
(Theorem~\ref{thm:shadow-galois-chenevier}) realise $p$-adic LOCAL
Galois representations; the $\MGSL$-action on the Universal Trace
Identity (Theorem~\ref{thm:mgsl-uti}) is compatible at the
period-ring level.

Both the chain-level LOCAL lane and the $(\infty, 1)$-categorical
GLOBAL lane are equally load-bearing: neither is a shadow of the
other (``lane assignment'' Pattern 269, \texttt{CLAUDE.md}).
\end{theorem}

\subsection*{The commuting cube: five theorems, two lanes, one moduli}

The master correspondence
Theorem~\ref{thm:master-LG} assembles into a single commuting cube:
the five theorems $\{A, B, C, D, H\}$ are the front face, the LOCAL
and GLOBAL forms occupy the top and bottom faces, and the
local--global maps are the vertical edges. The lateral edges between
theorems are the structural dependencies: Theorem $A$ feeds Theorem
$B$ (Positselski via bar-cobar), Theorem $C$ feeds Theorem $D$
(complementarity as Hodge supertrace feeding into Chern--Weil), Theorem
$H$ is transverse (concentration is an independent structural
feature).

Explicitly:
\[
\begin{array}{c@{\ }c@{\ }c@{\ }c@{\ }c@{\ }c@{\ }c}
   &            & A^{\mathrm{LOC}}   &            & B^{\mathrm{LOC}}   &            &                    \\
   & \swarrow   &                    & \searrow   &                    & \searrow   &                    \\
C^{\mathrm{LOC}}   &            & D^{\mathrm{LOC}}   &            & H^{\mathrm{LOC}}   &            &                    \\
   & \downarrow &                    & \downarrow &                    & \downarrow &                    \\
C^{\mathrm{GLOB}}  &            & D^{\mathrm{GLOB}}  &            & H^{\mathrm{GLOB}}  &            &                    \\
   & \nearrow   &                    & \nwarrow   &                    & \nwarrow   &                    \\
   &            & A^{\mathrm{GLOB}}  &            & B^{\mathrm{GLOB}}  &            &                    \\
\end{array}
\]

Commutativity witnesses: (i) $A$-row: fibrewise Koszul reflection
feeds Positselski (Thms~\ref{thm:A-local-global},
\ref{thm:B-local-global}); (ii) $C$-$D$ column: local-constancy
feeds Chern--Weil (Thms~\ref{thm:C-local-global}--\ref{thm:D-local-global});
(iii) $H$-column: fibrewise concentration extends to
$t$-exactness (Thm~\ref{thm:H-local-global}). The horizontal face is
Theorem~$A^{\infty, 2}$ (Francis--Gaitsgory,
\texttt{theorem\_A\_infinity\_2.tex}); the vertical face is
fibre-restriction functoriality of factorisation D-modules.

The cube is the Manin--Gaitsgory skeleton of Vol I: every theorem has
a position, every position has two faces, every face pair carries a
named morphism specified as a theorem rather than a convention.

\subsection*{Why the local--global discipline matters}

The discipline addresses three hidden confusion patterns.
\emph{(X1) Lane unnamed.} A chain-level proof on fixed $C$ does not
extend to $\Moduli_{g,n}$ for free; the chain-level $E_3$ identification
for affine KM at non-critical level
(\texttt{e3\_identification\_chain\_level\_platonic.tex}) is LOCAL, and
its GLOBAL lift via relative SC$^{\mathrm{ch,top}}$ is a separate,
not-yet-closed result.
\emph{(X2) Locally constant $\neq$ numerical.} Theorem $C$'s
$\kappa + \kappa^! \in \{0, 8, 13, 250/3, 98/3\}$ is a Hodge-rigidity
statement about locally constant sheaves on $M_{g,n}$, not a bare
number per genus.
\emph{(X3) Base-of-definition.} Chenevier determinants at $p = 691$
require $S = \Spec(\ZZ)$; $\MGSL$-action requires the rational motivic
base. Implicit base changes alter content silently.

Forward-looking inscriptions: LOCAL lane $=$ chain-level chart on
$\Conf_n(C)$ for fixed $C/\QQ$; GLOBAL lane $=$ fibred
$(\infty,1)$-category over $\Moduli_{g,n}/\Spec(\QQ)$; the
local--global theorem names the restriction functor and the
nodal-boundary behaviour explicitly.

\subsection*{Open residuals after the heal}

Three residuals remain after the five attack--heal cycles:

\begin{enumerate}
\item Theorem $B$ (Positselski) at the nodal stratum
$\partial\Moduli_{g,n}$: conjectural; the co-derived category may
require further resolution at the boundary. See Vol I
\texttt{theorem\_B\_scope\_platonic.tex} frontier list.
\item $F_5$ Drinfeld double at $g \ge 1$: universal Drinfeld
double over families; open even at $g = 1$ beyond Enriquez's
elliptic associator (genus-$1$ KZB in full generality pinned to
Drinfeld-double structure remains a programme frontier,
\texttt{feedback\_yangian\_type\_distinction.md}).
\item Deligne--Ihara injectivity: $\Gmot_{\MT(\ZZ)} \hookrightarrow
\GRT(\QQ)$ is surjective (Brown 2012); conjectured injective
(Deligne 2010). If non-injective, the cohomological-scope $\GRT(\QQ)$-
action factors through a proper quotient; the MGSL compatibility in
Theorem~\ref{thm:mgsl-uti} then carries a further torsor correction.
\end{enumerate}

None of these disturb the platonic form of the local--global
correspondence at the smooth-moduli open locus; all three are
named boundary / arithmetic frontier items, compatible with the Vol I
frontier list of §8 of the master synthesis.

\subsection*{Primary-source anchor list}

\begin{itemize}
\item Beilinson--Drinfeld 1999 \emph{Chiral Algebras} (AMS Colloq.\ 51,
2004 ed.): ch.~3 factorisation D-modules on Ran; \S3.4 relative
factorisation under nodal degeneration.
\item Francis--Gaitsgory 2012 ``Chiral Koszul duality,''
\emph{Selecta Math.}\ 18:27--87: $(\infty,1)$-categorical bar-cobar.
\item Gaitsgory 2007 ``Notes on geometric Langlands'' (extended
Whittaker): relative factorisation D-modules on moduli of curves.
\item Lurie 2017 \emph{Higher Algebra} §5.5: adjoint functor theorem
for presentable $(\infty,1)$-categories.
\item Drinfeld 1985 \emph{Soviet Math.\ Dokl.}\ 32 (Yangian/RTT);
Drinfeld 1989 \emph{Leningrad Math.\ J.}\ 1 (associator/pentagon).
\item Kohno 1987 \emph{Inv.\ Math.}\ 89: Drinfeld--Kohno algebra.
\item Enriquez 2008 arXiv:math/0607573 ($\GRT^{\mathrm{ell}}(\QQ)$);
Calaque--Enriquez--Etingof 2009 \emph{Math.\ Ann.}\ 345 (universal KZB).
\item Brown 2012 \emph{Ann.\ of Math.}\ 175 (arXiv:1102.1312):
$\Gmot_{\MT(\ZZ)} \twoheadrightarrow \GRT(\QQ)$.
\item Deligne--Goncharov 2005 \emph{Ann.\ Sci.\ ENS}\ 38; Willwacher
2015 \emph{Inv.\ Math.}\ 200 (arXiv:1009.1654); Tamarkin 1998
arXiv:math/9803025.
\item Kontsevich--Zagier 2001 ``Periods,'' in \emph{Mathematics
Unlimited 2001 and Beyond}: period ring $\Periods$.
\item Ihara 1990 ICM: Galois action on geometric $\pi_1$.
\item Kubota--Leopoldt 1964 \emph{J.\ reine angew.\ Math.}\ 214--215
($p$-adic $L$-functions, Kummer congruences); Ribet 1976 \emph{Inv.\
Math.}\ 34 (Herbrand--Ribet at $p = 691$).
\item Chenevier 2014 \emph{LMS LNS}\ 414: determinants in characteristic
$p \mid 2n!$.
\item Fontaine 1982 \emph{Ann.\ of Math.}\ 115 ($p$-adic Hodge); Tate
1967 \emph{Proc.\ Local Fields} ($p$-divisible groups).
\item Mumford 1977 \emph{Inv.\ Math.}\ 42: Hodge bundle Chern classes
on $\Moduli_{g,n}$.
\item Saito 1988/1990 \emph{Publ.\ RIMS Kyoto}\ 24/26: Hodge modules
on moduli, local constancy of supertraces.
\item Beilinson--Bernstein--Deligne 1982 \emph{Ast\'erisque}\ 100:
$t$-exactness.
\item Costello--Gwilliam 2021 \emph{Factorization Algebras in QFT}, Vol.\
2: super-BV, $\ZZ/2$-graded periods.
\end{itemize}

\subsection*{Cross-anchor discipline for forward inscriptions}

Every Vol I inscription henceforth: (1) specify LOCAL chart
(fixed $C/\QQ$, $g$, marked set) or state the theorem inherently
GLOBAL; (2) GLOBAL specifies base $\Spec(\QQ)/\Spec(\ZZ)$/etale;
(3) name the restriction functor and nodal-boundary behaviour;
(4) chain-level gives differential / homotopy / MC element / OPE pole,
$(\infty,1)$-categorical gives adjunction / $t$-structure /
limit-colimit construction. Already in use in
\texttt{mc5\_class\_m\_chain\_level\_platonic.tex} (chain-level LOCAL,
weight-completed) and \texttt{theorem\_A\_infinity\_2.tex}
($(\infty,2)$-GLOBAL Francis--Gaitsgory).

The Vol I climax chapter (nascent) states $d_{\mathrm{bar}} =
\KZ^*(\nabla_{\mathrm{Arnold}})$ and $K(A) = -c_{\mathrm{ghost}}(\mathrm{BRST}(A))$
at both lanes: LOCAL on $\Conf_n(C)$ (chain-level differential and
conductor) and GLOBAL on $\Moduli_{g,n}$ (flat connection, locally
constant ghost charge). The five theorems become corollaries of this
two-equation identity with their local--global refinements per
Theorem~\ref{thm:master-LG}.
